\section{Physical layer}
\label{sec:physical}

\subsection{Connection inherited from the DAQ--PDS ICD}

The physical CCM interface is taken from the \emph{Control, Configuration, and
Monitoring (CCM)} subsection of the DAQ--PDS ICD, EDMS 2088726, as maintained in
\code{daphne-icd} [R1, R2]:

\begin{quote}
Each DAPHNE and each LCM provides a dedicated 1\,GbE SFP link reserved for CCM
(no data readout on this link).

Fibers up to the PDS mini-rack patch panels are provided by PDS; continuation
and active switching are coordinated with DAQ and Technical Coordination.
\end{quote}

For DAPHNE, this means \textbf{one board Ethernet connection} shared by its
control, configuration and monitoring services. The physical path is:
\begin{center}
\textbf{DAPHNE Ethernet interface $\longleftrightarrow$ switch
$\longleftrightarrow$ physical server.}
\end{center}
The switch forwards Ethernet traffic to the appropriate server interface.
OPC UA, ZeroMQ, Protobuf and IPBus describe software communication carried by
the network. The OPC UA server is software deployed on physical server
infrastructure; it is not the name of a board connector or a separate cable.

\subsection{FD-VD inventory}

\begin{longtable}{L{0.48\textwidth}L{0.45\textwidth}}
\caption{FD-VD channel and board inventory.}\\
\toprule
\textbf{Item} & \textbf{FD-VD baseline} \\
\midrule
\endhead
PDS channels & 1,344 \\
Channels allocated to each DAPHNE & 32 \\
DAPHNE boards & $1,344/32=42$ \\
DAPHNE CCM connections & 42 board connections, one per board \\
Membrane PDS & 352 tiles, two channels per tile: 704 channels [R2] \\
Cathode PDS & 320 tiles, two channels per tile: 640 channels [R2] \\
\bottomrule
\end{longtable}

The membrane and cathode populations add to 1,344 channels. Membrane devices
use conductive power and electrical signal paths; cathode devices use
Power-over-Fiber and optical signal transmission [R2]. Both reach the same
DAPHNE CCM boundary. The installed mapping shall identify the population
served by each board. The 32-channel allocation is the FD-VD system mapping.

Servers, switches, PoF, LCM and power-supply units have separate installation
inventories.

\subsection{Physical topology and server boundary}

\begin{figure}[htbp]
\centering
\begin{tikzpicture}
\node[box,text width=30mm] (board) at (0,0)
 {DAPHNE\\CCM SFP\\1\,Gb Ethernet};
\node[box,text width=28mm] (switch) at (5.0,0)
 {Control-network\\switch};
\node[box,text width=36mm] (host) at (10.1,0)
 {Physical server\\network interface};
\draw[both] (board) -- node[above,yshift=6mm,align=center,font=\footnotesize]
 {Fibre via PDS\\patch panel} (switch);
\draw[both] (switch) -- node[above,yshift=6mm,align=center,font=\footnotesize]
 {Provisioned\\server link} (host);
\node[font=\footnotesize,align=center,text width=135mm] at (5.0,-1.65)
 {Repeated for 42 FD-VD DAPHNE board connections. Server and switch counts,
 uplinks and process placement are installation choices.};
\end{tikzpicture}
\caption{Physical CCM connection. Software-service placement is shown separately
in Figure~\ref{fig:protocol}.}
\end{figure}

The external OPC UA server hosts the integrated DAQ configuration component
and the Protobuf/ZeroMQ client of the DAPHNE service. Ignition connects as an
OPC UA client. The Hermes client remains separate from this OPC UA server.
It may run on the same physical host or another DAQ host; in either case its
traffic reaches the same board Ethernet interface. Software separation does
not require a second DAPHNE network connection.

SC provides computing infrastructure for its DCS processes.
PDS, SC and DAQ shall record the physical hosts for the
integrated OPC UA/configuration server and the separate Hermes client, with
installation and support assignments. The DPS controller and its physical I/O
require their own endpoint description.

\subsection{Physical interface specification}

\begingroup\small
\begin{longtable}{L{0.23\textwidth}L{0.38\textwidth}L{0.31\textwidth}}
\caption{CCM physical specification and remaining installation details.}\\
\toprule
\textbf{Element} & \textbf{Defined baseline} & \textbf{Installation detail to record} \\
\midrule
\endhead
DAPHNE port & One 1\,GbE SFP CCM link per board, copied from DAQ--PDS. & Board port label and installed SFP part number. \\
PDS fibre and patching & PDS supplies fibres to its mini-rack patch panels. & Fibre type, connector, wavelength, length, patch route and link budget. \\
Switch board-facing port & Compatible Ethernet termination for the board's 1\,GbE link. & Switch/port identity, compatible optic, agreed duplex and negotiation settings. \\
Switch uplinks & Carry the aggregate CCM traffic to the server network. & Topology, line rates, contention points and redundancy arrangement. \\
Physical server interface & Ethernet NIC connected to the provisioned control network. & Host/NIC identity, port, medium, rate and assignment to the server process. \\
\bottomrule
\end{longtable}
\endgroup

The cited CCM paragraph specifies an SFP link but does not supply the complete
CCM optical specification. The timing and 10\,GbE readout optical specifications
in DAQ--PDS apply to their respective links and shall not be copied into the
CCM link specification. Final installation drawings shall resolve the open
physical details above using the selected equipment.

The 1\,Gb/s line rate is shared by all enabled services on each board connection.
Uplink and server capacity shall be sized from the combined traffic and
concurrent operations in Section~\ref{sec:contention}.

DAQ acquisition data and the optical DTS signal retain the separate physical
interfaces defined by DAQ--PDS [R1, R2]. Here they are referenced only to identify
what their configuration and software services do over the shared CCM link.

Installation checks shall identify all 42 board routes and establish media
compatibility, link operation and simultaneous DAPHNE/Hermes reachability
from the actual external server infrastructure.
